\documentclass{article}
\usepackage[T1]{fontenc}
\usepackage{lmodern}
\usepackage{graphicx}
\usepackage{amsmath,amssymb}
\usepackage[a4paper,margin=2cm]{geometry}
\usepackage[absolute,overlay]{textpos}
\setlength{\TPHorizModule}{1mm}
\setlength{\TPVertModule}{1mm}
\usepackage[indonesian]{babel}
\usepackage{hyperref}
\setlength{\emergencystretch}{2em}
% unit: Halaman 89

\begin{document}

% === Halaman 89 (python/native) ===

89

def EkstraksiPesan(stego):

    audio = wave.open(stego, mode='rb')

    frames = bytearray(list(audio.readframes(audio.getnframes())))

    audio.close()

    \# Ambil semua LSB

    bit\_pesan = ''

    for byte in frames:

        bit\_pesan = bit\_pesan + str(byte \& 1)

   \# Kelompokkan setiap 8-bit dari bit\_pesan, simpan sebagai list          

    bit\_pesan = [bit\_pesan[i:i+8] for i in range(0, len(bit\_pesan), 8)]


    \# Ubah setiap 8-bit pesan menjadi sebuah karakter   

    pesan = ""    

    for i in range(len(bit\_pesan)):       

        if pesan[-5:] == "\#\#\#\#\#":   \# ketemu delimiter     

           break       

        else:            

           pesan = pesan + chr(int(bit\_pesan[i], 2))    

    \# Cari penanda akhir pesan

    end\_marker = pesan.find('\#\#\#\#\#')

    if end\_marker != -1:

        print("Pesan disisipkan:", pesan[:end\_marker])

    else:

        print("Tidak ditemukan pesan tersembunyi.")

\end{document}
